\documentclass[11pt]{article}

\usepackage[margin=1in]{geometry}
\usepackage{setspace}
\usepackage{amsmath, amssymb}
\usepackage{graphicx}
\usepackage{booktabs}
\usepackage{hyperref}

\title{Predicting Post-Graduation Earnings Using Machine Learning on U.S. College Scorecard Data}
\author{COMP 560 — Artificial Intelligence}
\date{}

\begin{document}

\maketitle

\onehalfspacing

\section*{Abstract}

We investigate whether institutional features such as tuition cost, school type, student demographics, and academic selectivity can predict median post-graduation earnings for U.S. college students. Using the U.S. Department of Education College Scorecard dataset (2023), comprising 950 institutions after preprocessing, we train and compare two models: a Linear Regression baseline and a Feedforward Neural Network (FFN) trained via backpropagation with dropout regularization. Our FFN achieves an $R^2$ of 0.626 on held-out test data, outperforming the linear baseline ($R^2 = 0.587$), demonstrating that salary outcomes exhibit non-linear structure. Feature analysis reveals that SAT math scores are the strongest predictor of graduate earnings—substantially stronger than tuition cost—suggesting academic selectivity matters more than price. To demonstrate practical applicability, we deploy the trained FFN behind an interactive Gradio interface that combines model predictions with LLM-generated (Claude) career guidance.

\section{Introduction}

The rising cost of higher education in the United States—now exceeding \$1.7 trillion in student loan debt—has intensified scrutiny over the return on investment of a college degree. Tuition continues to outpace wage growth, forcing students to make high-stakes financial decisions with limited accessible data. A central question emerges: does paying more for college actually result in higher earnings after graduation? This problem represents a natural application of machine learning, where institutional features can be used to predict long-term financial outcomes.

In this work, we construct a clean, feature-selected regression dataset from raw government data and compare linear and non-linear models for predicting post-graduation earnings. We evaluate a Linear Regression baseline against a Feedforward Neural Network (FFN), analyze feature importance, and identify which institutional characteristics most strongly influence earnings. Finally, we deploy an interactive tool that integrates ML predictions with LLM-generated career guidance, demonstrating an end-to-end application of course concepts including regression, neural networks, and transformers.

\section{Dataset}

\subsection{Data Source and Features}

The dataset is derived from the U.S. Department of Education College Scorecard (April 2023 release), which includes approximately 6,500 institutions and over 3,000 features. From this dataset, we select nine features relevant to earnings prediction. The target variable is \texttt{MD\_EARN\_WNE\_P10}, representing median earnings ten years after enrollment.

\begin{table}[h]
\centering
\caption{Selected Features and Descriptions}
\begin{tabular}{ll}
\toprule
Feature & Description \\
\midrule
NET\_PRICE & Net tuition cost \\
CONTROL & Public/Private indicator \\
UGDS & Undergraduate enrollment \\
PCTPELL & Pell Grant percentage \\
PCTFLOAN & Federal loan percentage \\
GRAD\_DEBT\_MDN\_SUPP & Median graduate debt \\
RPY\_3YR\_RT\_SUPP & Loan repayment rate \\
SATVRMID & SAT verbal midpoint \\
SATMTMID & SAT math midpoint \\
\bottomrule
\end{tabular}
\end{table}

\subsection{Preprocessing Pipeline}

Privacy-suppressed entries are replaced with NaN values. Public and private net price columns are merged into a unified \texttt{NET\_PRICE} feature. Rows containing missing values are removed, resulting in a final dataset of 950 institutions. Features are standardized using \texttt{StandardScaler}, fit exclusively on the training set to prevent data leakage.

\subsection{Dataset Limitations}

The dataset includes only students receiving federal financial aid, excluding those who pay full tuition without aid. This introduces selection bias, particularly at elite institutions where high-income students are underrepresented. As a result, the model may underestimate true earnings at highly selective universities.

\section{Methods}

\subsection{Problem Formulation}

We define the regression problem as learning a function $f: \mathbb{R}^9 \rightarrow \mathbb{R}$ such that $f(x_i) \approx y_i$. The dataset is split into 80\% training and 20\% testing, with five-fold cross-validation performed on the training set.

\subsection{Linear Regression Baseline}

The baseline model minimizes Mean Squared Error (MSE):

\[
L(w) = \frac{1}{N} \sum_{i=1}^{N} (y_i - w^T x_i - b)^2
\]

Linear Regression provides interpretability and serves as a strong baseline without requiring hyperparameter tuning.

\subsection{Feedforward Neural Network}

The FFN architecture is defined as:

\begin{center}
Input (9) $\rightarrow$ 128 $\rightarrow$ 64 $\rightarrow$ 32 $\rightarrow$ 1
\end{center}

ReLU activations and dropout (0.2) are applied between layers. The model is trained using the Adam optimizer (learning rate $10^{-3}$), with MSE loss, batch size 32, and 150 epochs.

\subsection{Reproducibility and Model Persistence}

Random seeds are fixed across all libraries. Models are saved as \texttt{ffn\_model.pt}, \texttt{lr\_model.pkl}, and \texttt{scaler.pkl}. The codebase supports both local and cloud execution environments.

\subsection{Interactive Deployment}

The trained FFN is deployed using a Gradio interface and integrated with the Claude API for natural language explanations. The pipeline maps user input to predictions and generates contextualized career recommendations.

\section{Results}

\subsection{Quantitative Model Comparison}

\begin{table}[h]
\centering
\caption{Model Performance Comparison}
\begin{tabular}{lccc}
\toprule
Model & MAE (\$) & RMSE (\$) & $R^2$ \\
\midrule
Linear Regression & 6596 & 8512 & 0.587 \\
FFN & 6148 & 8101 & 0.626 \\
\bottomrule
\end{tabular}
\end{table}

The FFN outperforms Linear Regression across all metrics.

\subsection{FFN Training Dynamics}

\begin{figure}[h]
\centering
\includegraphics[width=0.6\textwidth]{loss_curve.png}
\caption{Training Loss Over Epochs}
\end{figure}

Loss decreases steadily, indicating stable convergence.

\subsection{Feature Importance}

\begin{table}[h]
\centering
\caption{Feature Importance (Linear Regression Coefficients)}
\begin{tabular}{lc}
\toprule
Feature & Effect \\
\midrule
SAT Math & +10142 \\
Loan Repayment Rate & +4878 \\
Net Price & +2359 \\
SAT Verbal & -2867 \\
\bottomrule
\end{tabular}
\end{table}

\section{Discussion}

SAT scores emerge as a significantly stronger predictor of earnings than tuition cost, supporting the hypothesis that student quality outweighs institutional price. The FFN achieves modest improvement over Linear Regression, suggesting mostly linear relationships with some non-linear interactions. Limitations include dataset bias and lack of major-level features. The interactive tool demonstrates practical deployment of ML models.

\section{Conclusion}

The FFN outperforms Linear Regression ($R^2 = 0.626$ vs. $0.587$), confirming the presence of non-linear structure in earnings prediction. SAT scores are the strongest predictor, indicating that academic selectivity is more informative than tuition cost. The deployed interactive system highlights real-world applicability.

\section{Future Work}

Future work includes incorporating multi-year data, adding major-level features, addressing selection bias with external datasets, and exploring ensemble methods.

\section{References}

\begin{itemize}
\item U.S. Department of Education. College Scorecard Data. 2023. \url{https://collegescorecard.ed.gov/data}
\item Chetty et al. Mobility Report Cards. NBER, 2017.
\item Goodfellow et al. Deep Learning. MIT Press, 2016.
\item Anthropic. Claude Documentation. 2025.
\item U.S. Department of Education. Technical Documentation. 2023.
\end{itemize}

\end{document}